% Appendix A: material compressed out of Section 2 on Day 4 (2026-09-12).
% Moved verbatim; nothing was deleted from the paper, only relocated.
\section{Testbed, harness and audit protocol: full detail}
\label{app:testbed}

\subsection{Training harness, policy classes, and measurement}

All experiments run in one single-file PPO harness {[}PPO; GAE{]}
differing only in the policy head, with GAE(\(\lambda\)) advantages,
clipped surrogate (clip 0.2, 4 epochs, 4 minibatches), Adam at
\(3 \times 10^{-4}\) annealed linearly to zero, gradient-norm cap 0.5,
16 parallel environments, 128-step rollouts, and 15M environment steps
per run; the harness discount equals the model's \(\beta\) per instance.
Policy classes: \textbf{mlp} (two Tanh layers of 128), \textbf{gru128}
(GRU with 128 units) as the memory arm on external benchmarks, and two
restricted heads that exist for identification rather than performance:
\textbf{linear} (\(u = -kv\), a single scalar read from the raw
observation, no bias, no trunk) and \textbf{linearb} (\(u = -kv + b\)).
The linear head cannot leave the gain manifold, so any instability it
fails to show is thereby located off that manifold; the linearb head
adds exactly one escape direction, the one whose stability the theory
computes. Both restricted heads support initialisation at a chosen gain,
and linearb supports kicking \(b\) to a chosen \(b_0\) at \(t=0\),
turning stability claims into perturbation-decay measurements rather
than waits for stochastic escape. All new parameters are
constant-initialised, and we verified the harness produces
bitwise-identical mlp and gru128 trajectories, rows and curves, before
and after these additions: the registered E6/E7/E9 data remains
reproducible from the released file.

\textbf{Deployed evaluation.} Reported cost is measured with the policy
mean, no exploration noise, on a single 200\{,\}000-step trajectory with
a 1\{,\}000-step burn-in discarded, giving the stationary cost rate
\(J_{\text{rate}}\). The obvious alternative --- averaging many short
episodes --- is wrong here in a way that would have manufactured a
registered result: episodes reset from the uncontrolled stationary
distribution, every episode pays the settling transient, and the
inflation is multiplicative and grows with the gain (measured factors
1.088, 1.270, 2.415 at \(k\) = 1.65, 5.99, 20.55), i.e.~it is largest
exactly where the predicted effect is largest and \(\lambda\) moves the
gain. A calibration test pins policies at known gains and requires the
measured-to-analytic ratio to be flat in \(k\); it passes with spread
\(1.6\times10^{-4}\) across gains, which is what licenses comparing
\(J_{\text{rate}}\) to the frozen baselines at all. Two diagnostics are
read from every continuous run: the \textbf{fitted gain} (the OLS slope
of \(u\) on \(v\) over the evaluation trajectory) and, for linearb, the
terminal bias \(b\). Training-time behaviour is summarised by the
\textbf{intermittency fraction}: the fraction of late-half evaluation
points whose cost exceeds \(2 J_{\text{ml}}\), which detects the
recurrent-excursion phenomenology (section 5) without reference to any
equilibrium prediction.

\textbf{External benchmarks.} The corroboration arm (E9) uses POPGym
{[}POPGym{]} environments reimplemented natively --- RepeatPrevious and
position-only CartPole --- with the upstream seeding path reproduced
exactly (PCG64 via the same SeedSequence construction) and verified
bitwise against the reference implementation over 13\{,\}542 steps
before any run. This removes the cluster dependency chain entirely; jobs
run on plain numpy/torch with no per-node compilation.


\subsection{Audit protocol}

The claims are adjudicated by preregistration, and the preregistrations
are \textbf{executable}: frozen constants and scoring functions live in
Python files whose freeze is asserted mechanically --- the deployment
script runs every prereg as a pre-flight and refuses to submit if any
sentinel remains --- so a prediction cannot drift between freeze and
verdict. Amendments are dated, and are made before the relevant data
exists except where explicitly disclosed otherwise in the amendment text
itself. Scorers for the experiments whose results files accumulate
across submissions (E6/E7/E9, E12, E13) refuse to run on data with
duplicated cells; the E12 guard was added post-data after its grid
re-ran the probe's two reference cells, an episode \S6 reports along with
the verdict-neutral separation. Two departures from the discipline are
disclosed rather than smoothed over, because a pre-submission audit
found them and an audit paper does not get to hide its own findings: the
scoring code for P10.1--P10.3 was implemented after the data existed
(the predictions and constants were frozen on 2026-08-14; the file's
advertised \texttt{-\/-score} mode was a stub until 2026-08-21, and the
implemented scorer prints a sensitivity table over the one aggregation
choice its freeze had left unspecified), and the diagnostic prereg (E8)
lacked a mechanical freeze assertion until the same date. Verdicts are
published as registered: this paper reports {[}N{]} registered failures
alongside its passes, and two of its analyses (sections 6.x, 6.y) are
labelled post-hoc and carry no confirmatory weight. Because cluster
nodes of different CPU generations produce bit-different trajectories
from identical seeds, the node class of every job is recorded and
reported; within-class runs reproduce bitwise. Across classes we
measured seven perturbed reruns: the largest difference is about 1.9
across-seed standard deviations, the typical one far below one, and
substituting the affected values changes no registered verdict
{[}APPENDIX CROSS-REF{]}. One comparison inherits a hardware caveat from
this: E9 ran entirely within the d class while E6 and E7 both include
t-class tasks, so P7.2 --- an ordering test across instances ---
compares cells that were not all produced on bitwise-comparable
hardware. P7.2 failed as registered, so no claim rests on it; the
heterogeneity is recorded because it is part of what that comparison
was. Seeds are 1--5 for GLEAM experiments (1--3 for E10 and E13),
matched across arms and architectures so that comparisons are paired at
the seed level.


\subsection{Number provenance}
\label{app:provenance}

\begin{longtable}[]{@{}
  >{\raggedright\arraybackslash}p{(\columnwidth - 2\tabcolsep) * \real{0.5000}}
  >{\raggedright\arraybackslash}p{(\columnwidth - 2\tabcolsep) * \real{0.5000}}@{}}
\toprule\noalign{}
\begin{minipage}[b]{\linewidth}\raggedright
constant in text
\end{minipage} & \begin{minipage}[b]{\linewidth}\raggedright
generated by
\end{minipage} \\
\midrule\noalign{}
\endhead
\bottomrule\noalign{}
\endlastfoot
\(F, G, \Sigma_d\), \(J(k)\), \(k_{\text{ml}}\), \(J_{\text{ml}}\),
\(J^\ast\) & \texttt{vendor/gleam/theorem2\_learning.py}
(\texttt{disc\_model}, \texttt{J\_deploy}, \texttt{best\_ml},
\texttt{J\_opt}) \\
\(k_{\text{AC}}(\lambda)\), LSTD(\(\lambda\)) fixed point, resolvent
identities & \texttt{vendor/gleam/lstd\_lambda.py}
(\texttt{k\_eq\_lambda}, \texttt{lstd\_lambda\_theta2}) \\
\(\lambda_{1/2}\) per instance, grid-matched targets,
\(J(k_{\text{AC}})\) tables & \texttt{prereg\_phase2.py} (REFERENCE
block), \texttt{prereg\_e10.py} (\texttt{C}) \\
\(S(\lambda,k)\), \(\lambda^\ast\) per instance &
\texttt{vendor/gleam/bias\_stability.py} (\texttt{bias\_stability},
\texttt{predict\_all}), frozen in \texttt{prereg\_e7\_bias.py} \\
eval inflation 1.088/1.270/2.415, calibration spread
\(1.6\times10^{-4}\) & \texttt{harness/test\_eval\_calibration.py};
prereg\_phase2.py header note (dated 2026-08-12) \\
bitwise env verification, 13,542 steps &
\texttt{harness/test\_env\_equivalence.py} \\
bitwise harness regression after linear/linearb & MEMO-arch-linear.md \S6
(rows + curves, timing columns excluded) \\
\(k_{\text{AC}}(0)=24.79\) on g1\_h0.05 (the TD(0) value; constant in
every raw row as \texttt{ref\_k\_AC}) &
\texttt{vendor/gleam/gleam\_bench.py} (\texttt{ac\_equilibrium}, Props
1--3); reproduced by \texttt{lstd\_lambda.py} at \(\lambda=0\) \\
\(k_{\text{AC}}(0.9)=20.55\), \(k_{\text{ml}}=1.649\) on g1\_h0.05 &
\texttt{prereg\_phase2.py} \texttt{K\_AC\_LAMBDA}, \texttt{BASELINES} \\
intermittency definition (\(2J_{\text{ml}}\), late half, 0.05 flag) &
\texttt{a5-tier1b-findings.md} \S1 (frozen 2026-08-13); constants carried
in \texttt{prereg\_e7\_bias.py} (\texttt{INTERMITTENCY\_*},
\texttt{LATE\_FRACTION}) \\
\end{longtable}

